\documentclass[11pt]{article}

\usepackage[margin=1in]{geometry}
\usepackage{amsmath,amssymb}
\usepackage{tikz}

\title{FA2}
\author{Awit, Julia Theressa}
\date{August 29, 2026}

\begin{document}
\maketitle

\section*{Problem 3.1}

Given the data matrix
\[
X=\begin{pmatrix} 9 & 1 \\ 5 & 3 \\ 1 & 2\end{pmatrix}.
\]

First, find the sample mean of each column.
\[
\bar x_1=\frac{9+5+1}{3}=5, \qquad \bar x_2=\frac{1+3+2}{3}=2
\]
\[
\bar{\mathbf x}=\begin{pmatrix}5\\2\end{pmatrix}
\]

\textbf{(a)} The three points are $(9,1)$, $(5,3)$, $(1,2)$. The sample mean $(5,2)$ is
plotted as the center of the three points.

\begin{center}
\begin{tikzpicture}[scale=0.5]
\draw[->] (-1,0) -- (11,0) node[right]{$x_1$};
\draw[->] (0,-1) -- (0,5) node[above]{$x_2$};
\foreach \x in {1,3,5,7,9} \draw (\x,-0.1)--(\x,0.1) node[below=3pt]{\scriptsize \x};
\foreach \y in {1,2,3,4} \draw (-0.1,\y)--(0.1,\y) node[left=3pt]{\scriptsize \y};
\filldraw (9,1) circle (2.5pt) node[above right]{$(9,1)$};
\filldraw (5,3) circle (2.5pt) node[above right]{$(5,3)$};
\filldraw (1,2) circle (2.5pt) node[above right]{$(1,2)$};
\filldraw (5,2) circle (2.5pt) node[above right]{$\bar{\mathbf x}=(5,2)$};
\end{tikzpicture}
\end{center}

\textbf{(b)} Now treat each column as a vector in $n=3$ dimensions.
\[
\mathbf y_1=\begin{pmatrix}9\\5\\1\end{pmatrix},\qquad
\mathbf y_2=\begin{pmatrix}1\\3\\2\end{pmatrix}
\]

Subtract the mean (times $\mathbf 1$) from each vector to get the deviation vectors.
\[
\mathbf y_1-\bar x_1\mathbf 1=\begin{pmatrix}9\\5\\1\end{pmatrix}-\begin{pmatrix}5\\5\\5\end{pmatrix}
=\begin{pmatrix}4\\0\\-4\end{pmatrix}
\]
\[
\mathbf y_2-\bar x_2\mathbf 1=\begin{pmatrix}1\\3\\2\end{pmatrix}-\begin{pmatrix}2\\2\\2\end{pmatrix}
=\begin{pmatrix}-1\\1\\0\end{pmatrix}
\]

These are the deviation vectors $\mathbf d_1=(4,0,-4)'$ and $\mathbf d_2=(-1,1,0)'$. In the
sketch, $\mathbf y_1$ and $\mathbf y_2$ are vectors in $\mathbb R^3$, and $\mathbf d_1$,
$\mathbf d_2$ are what is left after removing their projection onto the line through
$\mathbf 1=(1,1,1)'$.

\textbf{(c)} Now plot $\mathbf d_1$ and $\mathbf d_2$ starting from the origin.

\begin{center}
\begin{tikzpicture}[scale=1]
\draw[->] (-1.2,0) -- (4.5,0);
\draw[->] (0,-0.3) -- (0,2.5);
\draw[thick,->] (0,0) -- (3.5,0) node[right]{$\mathbf d_1$};
\draw[thick,->] (0,0) -- (-0.9,1.2) node[above]{$\mathbf d_2$};
\draw (0.5,0) arc (0:127:0.5);
\node at (0.05,0.55) {\scriptsize $\theta$};
\end{tikzpicture}
\end{center}

Length of $\mathbf d_1$:
\[
\lVert \mathbf d_1\rVert=\sqrt{4^2+0^2+(-4)^2}=\sqrt{32}=\boxed{4\sqrt2}
\]

Length of $\mathbf d_2$:
\[
\lVert \mathbf d_2\rVert=\sqrt{(-1)^2+1^2+0^2}=\boxed{\sqrt2}
\]

Dot product:
\[
\mathbf d_1\cdot \mathbf d_2=(4)(-1)+(0)(1)+(-4)(0)=-4
\]

Cosine of the angle between them:
\[
\cos\theta=\frac{\mathbf d_1\cdot \mathbf d_2}{\lVert \mathbf d_1\rVert\lVert \mathbf d_2\rVert}
=\frac{-4}{(4\sqrt2)(\sqrt2)}=\frac{-4}{8}=\boxed{-\frac12}
\]

Relating this to $S_n$: using $S_n=\frac1n D'D$,
\[
s_{11}=\frac{\lVert\mathbf d_1\rVert^2}{3}=\frac{32}{3},\qquad
s_{22}=\frac{\lVert\mathbf d_2\rVert^2}{3}=\frac23,\qquad
s_{12}=\frac{\mathbf d_1\cdot \mathbf d_2}{3}=-\frac43
\]
\[
S_n=\boxed{\begin{pmatrix} 32/3 & -4/3\\ -4/3 & 2/3\end{pmatrix}}
\]

The correlation is
\[
r_{12}=\frac{s_{12}}{\sqrt{s_{11}s_{22}}}=\frac{-4/3}{\sqrt{(32/3)(2/3)}}=\frac{-4/3}{8/3}=\boxed{-\frac12}
\]
which is exactly $\cos\theta$ from above. So $R=\begin{pmatrix}1 & -1/2\\ -1/2 & 1\end{pmatrix}$.
This shows that the length of a deviation vector is tied to the variance of that variable,
and the angle between two deviation vectors is tied to their correlation.

\newpage
\section*{Problem 3.6}

\[
X=\begin{pmatrix} -1 & 3 & -2\\ 2 & 4 & 2\\ 5 & 2 & 3\end{pmatrix}
\]

Column means:
\[
\bar x_1=\frac{-1+2+5}{3}=2,\qquad \bar x_2=\frac{3+4+2}{3}=3,\qquad \bar x_3=\frac{-2+2+3}{3}=1
\]

\textbf{(a)} Subtract the mean from each column to get $X-\mathbf 1\bar{\mathbf x}'$:
\[
X-\mathbf 1\bar{\mathbf x}'=\begin{pmatrix} -3 & 0 & -3\\ 0 & 1 & 1\\ 3 & -1 & 2\end{pmatrix}=D
\]

Is $D$ full rank? Add up the rows:
\[
(-3+0+3,\ 0+1-1,\ -3+1+2)=(0,0,0)
\]
The rows add to the zero vector, so the rows are linearly dependent, and $\operatorname{rank}(D)\le 2$.
Rows 1 and 2, $(-3,0,-3)$ and $(0,1,1)$, are not multiples of each other, so
$\operatorname{rank}(D)=2$. Since $D$ is $3\times3$, it is \textbf{not} full rank (full
rank would be $3$). This always happens for a mean-centered matrix because the columns of
$D$ each sum to zero.

\textbf{(b)} $S=\frac1{n-1}D'D=\frac12 D'D$.
\[
D'D=\begin{pmatrix}
18 & -3 & 15\\
-3 & 2 & -1\\
15 & -1 & 14
\end{pmatrix}
\]
\[
S=\frac12\begin{pmatrix}
18 & -3 & 15\\
-3 & 2 & -1\\
15 & -1 & 14
\end{pmatrix}
=\boxed{\begin{pmatrix}
9 & -3/2 & 15/2\\
-3/2 & 1 & -1/2\\
15/2 & -1/2 & 7
\end{pmatrix}}
\]

Since $D$ is not full rank, $D'D$ is singular, so
\[
|S|=\boxed{0}.
\]
Geometrically, $|S|$ measures the volume of the box formed by the deviation vectors. A
value of $0$ means the deviation vectors do not really span all $3$ dimensions --- they
are squashed into a flat (2-dimensional) region, so there is no real 3-D spread in the data.

\textbf{(c)} Total sample variance $=\operatorname{tr}(S)=9+1+7=\boxed{17}$.

\newpage
\section*{Problem 3.9}

\[
X=\begin{pmatrix}
12 & 17 & 29\\
18 & 20 & 38\\
14 & 16 & 30\\
20 & 18 & 38\\
16 & 19 & 35
\end{pmatrix}
\]
where $x_3=x_1+x_2$ (total score).

Column means:
\[
\bar x_1=16,\qquad \bar x_2=18,\qquad \bar x_3=34
\]
(and $34=16+18$, as expected.)

\textbf{(a)} Mean-corrected data matrix:
\[
D=X-\mathbf 1\bar{\mathbf x}'=
\begin{pmatrix}
-4 & -1 & -5\\
2 & 2 & 4\\
-2 & -2 & -4\\
4 & 0 & 4\\
0 & 1 & 1
\end{pmatrix}
\]

In every row, column 3 is column 1 plus column 2 (e.g. row 1: $-4+(-1)=-5$). So
\[
\text{col}_1+\text{col}_2-\text{col}_3=\mathbf 0
\]
which means
\[
\boxed{\mathbf a'=[a_1,a_2,a_3]=[1,\ 1,\ -1]}
\]
gives the linear dependence.

\textbf{(b)} $S=\frac1{n-1}D'D=\frac14 D'D$.
\[
D'D=\begin{pmatrix}
40 & 12 & 52\\
12 & 10 & 22\\
52 & 22 & 74
\end{pmatrix}
\]
\[
S=\boxed{\begin{pmatrix}
10 & 3 & 13\\
3 & 5/2 & 11/2\\
13 & 11/2 & 37/2
\end{pmatrix}}
\]

Since the columns of $D$ are linearly dependent, $D'D$ (and $S$) is singular, so
$|S|=\boxed{0}$.

Check $S\mathbf a=\mathbf 0$ with $\mathbf a=(1,1,-1)'$:
\[
S\mathbf a=
\begin{pmatrix}
10+3-13\\
3+2.5-5.5\\
13+5.5-18.5
\end{pmatrix}
=\begin{pmatrix}0\\0\\0\end{pmatrix}
\]
So $\mathbf a$ is an eigenvector of $S$ with eigenvalue $0$.

\textbf{(c)} Checking the original data: $12+17=29$, $18+20=38$, $14+16=30$, $20+18=38$,
$16+19=35$. In each row $x_3=x_1+x_2$, confirming
\[
1\cdot x_1+1\cdot x_2+(-1)\cdot x_3=0,
\]
i.e. $a_1=1,\ a_2=1,\ a_3=-1$.

\newpage
\section*{Problem 3.18}

\[
\bar{\mathbf x}=\begin{pmatrix}0.766\\0.508\\0.438\\0.161\end{pmatrix},\qquad
S=\begin{pmatrix}
0.856 & 0.635 & 0.173 & 0.096\\
0.635 & 0.568 & 0.127 & 0.067\\
0.173 & 0.127 & 0.171 & 0.039\\
0.096 & 0.067 & 0.039 & 0.043
\end{pmatrix}
\]
($x_1$ petroleum, $x_2$ natural gas, $x_3$ hydroelectric, $x_4$ nuclear.)

\textbf{(a)} Total consumption is $T=x_1+x_2+x_3+x_4$, so $\mathbf a=(1,1,1,1)'$.

Mean:
\[
E[T]=0.766+0.508+0.438+0.161=\boxed{1.873}
\]

Variance ($=$ sum of all entries of $S$):
\[
\operatorname{Var}(T)=\mathbf a'S\mathbf a
\]
\[
=(0.856+0.635+0.173+0.096)+(0.635+0.568+0.127+0.067)+(0.173+0.127+0.171+0.039)+(0.096+0.067+0.039+0.043)
\]
\[
=1.760+1.397+0.510+0.245=\boxed{3.912}
\]

\textbf{(b)} Excess of petroleum over gas: $x_1-x_2$, so $\mathbf b=(1,-1,0,0)'$.

Mean:
\[
E[x_1-x_2]=0.766-0.508=\boxed{0.258}
\]

Variance:
\[
\operatorname{Var}(x_1-x_2)=s_{11}+s_{22}-2s_{12}=0.856+0.568-2(0.635)=\boxed{0.154}
\]

Covariance with the total $T$: $\mathbf b'S\mathbf a$. First find $\mathbf b'S$ (row 1 of
$S$ minus row 2 of $S$):
\[
(0.856-0.635,\ 0.635-0.568,\ 0.173-0.127,\ 0.096-0.067)=(0.221,\ 0.067,\ 0.046,\ 0.029)
\]
Then dot this with $\mathbf a=(1,1,1,1)'$:
\[
\operatorname{Cov}(x_1-x_2,T)=0.221+0.067+0.046+0.029=\boxed{0.363}
\]

\newpage
\section*{Problem 3.20}

\begin{center}
\begin{tabular}{cc|cc|cc}
$x_1$ & $x_2$ & $x_1$ & $x_2$ & $x_1$ & $x_2$\\ \hline
12.5 & 13.7 & 9.0 & 24.4 & 3.5 & 26.1\\
14.5 & 16.5 & 6.5 & 18.2 & 8.0 & 14.5\\
8.0 & 17.4 & 10.5 & 22.0 & 17.5 & 42.3\\
9.0 & 11.0 & 10.0 & 32.5 & 10.5 & 17.5\\
19.5 & 23.6 & 4.5 & 18.7 & 12.0 & 21.8\\
8.0 & 13.2 & 7.0 & 15.8 & 6.0 & 10.4\\
9.0 & 32.1 & 8.5 & 15.6 & 13.0 & 25.6\\
7.0 & 12.3 & 6.5 & 12.0 & & \\
7.0 & 11.8 & 8.0 & 12.8 & &
\end{tabular}
\end{center}
($x_1=$ storm duration, $x_2=$ hours to clear snow, $n=25$.)

\textbf{(a)} Using the summary statistics:
\[
\bar x_1=\frac{235.5}{25}=9.42,\qquad \bar x_2=\frac{481.8}{25}=19.272
\]
\[
\sum(x_{j1}-\bar x_1)^2=339.34\ \Rightarrow\ s_{11}=\frac{339.34}{24}=14.1392
\]
\[
\sum(x_{j2}-\bar x_2)^2=1493.7304\ \Rightarrow\ s_{22}=\frac{1493.7304}{24}=62.2388
\]
\[
\sum(x_{j1}-\bar x_1)(x_{j2}-\bar x_2)=323.344\ \Rightarrow\ s_{12}=\frac{323.344}{24}=13.4727
\]

Let $d=x_2-x_1$. Then
\[
\bar d=\bar x_2-\bar x_1=19.272-9.42=\boxed{9.852}
\]
\[
s_d^2=s_{11}+s_{22}-2s_{12}=14.1392+62.2388-2(13.4727)=\boxed{49.4326}
\]

\textbf{(b)} Now compute the differences directly, $d_j=x_{j2}-x_{j1}$:
\[
1.2,\ 2.0,\ 9.4,\ 2.0,\ 4.1,\ 5.2,\ 23.1,\ 5.3,\ 4.8,\ 15.4,\ 11.7,\ 11.5,\ 22.5,
\]
\[
14.2,\ 8.8,\ 7.1,\ 5.5,\ 4.8,\ 22.6,\ 6.5,\ 24.8,\ 7.0,\ 9.8,\ 4.4,\ 12.6
\]

Sum $=246.3$, so
\[
\bar d=\frac{246.3}{25}=\boxed{9.852}
\]
which matches part (a) exactly. The sum of squared deviations is $1186.367$, so
\[
s_d^2=\frac{1186.367}{24}=\boxed{49.4320}
\]

Both methods give the same mean ($9.852$) and basically the same variance ($49.4326$ vs
$49.4320$, the small difference is just rounding). This makes sense because
$\operatorname{Var}(x_2-x_1)=s_{11}+s_{22}-2s_{12}$ should hold either way, whether we
compute it from the summary statistics or straight from the differences.

\end{document}